\begin{figure}[ht]
\centering
\begin{tikzpicture}[x=1cm, y=1cm, font=\small]

% ---- Axes ----
\draw[->, thick] (0, 0) -- (10.6, 0) node[right, font=\footnotesize]
  {time $t / \tau$};
\draw[->, thick] (0, 0) -- (0, 5.4) node[above, font=\footnotesize]
  {response $y(t) / y_{\infty}$};

% Final-value asymptote at y_inf (mapped to y = 4.5)
\draw[thick, enggray, dashed] (0, 4.5) -- (10, 4.5);
\node[font=\scriptsize, color=enggray, anchor=south east] at (10, 4.5)
  {final value $y_{\infty}$ ($100\%$)};

% x-axis ticks in units of tau (2 cm per tau)
\foreach \x/\lab in {2/1, 4/2, 6/3, 8/4, 10/5} {
  \draw (\x, -0.1) -- (\x, 0.1);
  \node[font=\scriptsize, anchor=north] at (\x, -0.1) {\lab};
}
% y-axis gridlines at 63.2% and 90%
\draw[thin, engred, dotted] (0, 2.844) -- (2, 2.844) -- (2, 0);
\node[font=\scriptsize, color=engred, anchor=east] at (-0.1, 2.844)
  {$63.2\%$};
\draw[thin, engblue, dotted] (0, 4.05) -- (4.605, 4.05) -- (4.605, 0);
\node[font=\scriptsize, color=engblue, anchor=east] at (-0.1, 4.05)
  {$90\%$};

% First-order rising exponential y = 4.5 (1 - e^{-t/tau}); sampled points
% t/tau : y/y_inf : plotted (x=2*t/tau, yy=4.5*value)
% 0:0    0.5:0.393  1:0.632  1.5:0.777  2:0.865  3:0.950  4:0.982  5:0.993
\draw[very thick, engred]
  (0, 0)
  .. controls (0.5, 1.05) and (1.1, 1.95) ..
  (2, 2.844)
  .. controls (2.8, 3.45) and (3.6, 3.80) ..
  (4.605, 4.05)
  .. controls (5.6, 4.28) and (7.0, 4.42) ..
  (10, 4.485);

% mark tau63 point
\fill[engred] (2, 2.844) circle (1.8pt);
\node[font=\scriptsize, color=engred, anchor=west] at (2.1, 2.55)
  {$\tau_{63}$ at $t = \tau$};

% mark tau90 point
\fill[engblue] (4.605, 4.05) circle (1.8pt);
\node[font=\scriptsize, color=engblue, anchor=west] at (4.7, 3.75)
  {$\tau_{90}$ at $t \approx 2.30\,\tau$};

% initial-slope tangent line: slope y_inf/tau -> from (0,0) reaches 4.5 at t/tau=1 (x=2)
\draw[thin, engamber] (0, 0) -- (2.0, 4.5);
\node[font=\scriptsize, color=engamber, anchor=south west, rotate=0] at (0.15, 4.55)
  {initial-slope tangent (crosses $y_{\infty}$ at $t = \tau$)};

\end{tikzpicture}
\caption[First-order sensor step response]{The step response of a
first-order sensor, $y(t) = y_{\infty}(1 - e^{-t/\tau})$, normalised
to its final value. The sensor reaches $63.2\%$ of the step at one
time constant ($\tau_{63} = \tau$) and $90\%$ at $t \approx 2.30\,\tau
= \tau_{90}$. The initial-slope tangent (amber) crosses the final
value at exactly $t = \tau$, the graphical construction that recovers
the time constant from a measured trace. A datasheet that quotes a
$\tau_{63}$ response time and a datasheet that quotes a $\tau_{90}$
response time for the same physical sensor differ by the factor
$\ln 10 \approx 2.30$; reading which fraction the manufacturer used
is part of reading the datasheet.}
\label{fig:vol01:ch05:step-response}
\end{figure}
